 
\chapter{Simulation Results and Performance Benchmarking} \label{chap:Benchmarking}


\section{Simulation Setup and Input Test Signals}
\label{sec:simulation_setup}

To validate the theoretical framework established in Chapter 3 and evaluate the physical performance of the proposed asynchronous Flash ADC, extensive transistor-level simulations were conducted using a 16-nm CMOS process node under a 0.8~V nominal supply voltage. This section details the simulation environment, the specialized reconstruction methodology required for non-uniform asynchronous data, and the characteristics of the input test stimuli.

Unlike conventional synchronous data converters where output samples are bound to a periodic clock grid, the asynchronous level-crossing Flash ADC generates discrete output tokens only when a quantization threshold is traversed. Consequently, traditional Fast Fourier Transform (FFT) algorithms, which strictly require uniformly sampled time-domain records, cannot be directly applied to the raw asynchronous output. To accurately characterize the dynamic performance metrics such as Signal-to-Noise Ratio (SNR), Signal-to-Noise-and-Distortion Ratio (SNDR), and Effective Number of Bits (ENOB) a non-uniform to uniform time-domain resampling and mathematical reconstruction algorithm was implemented.

The primary input test signal employed for dynamic characterization is a full-scale sinusoidal waveform near the Nyquist limit, superimposed with transient thermal noise to model realistic operating conditions. The spectrum of the reconstructed asynchronous output, showcasing the fundamental signal bin, the noise floor, and harmonic distortion components, is illustrated in the FFT spectrum below.

\begin{figure}[htbp]
    \centering
    % \includegraphics[width=0.75\textwidth]{figures/fft_spectrum.pdf}
    \vspace{5cm} % Placeholder space for the image
    \caption{FFT spectrum of the reconstructed asynchronous Flash ADC output signal.}
    \label{fig:fft_spectrum}
\end{figure}

The dynamic performance curves, mapping the behavior of both SNDR and ENOB across a broad range of input frequencies and amplitudes, are provided in Figure~\ref{fig:dynamic_curves}. These curves contrast the ideal noiseless performance against the degraded parameters under transient noise conditions, tracking the incidence of noise-induced threshold chatter.

\begin{figure}[htbp]
    \centering
    % \includegraphics[width=0.75\textwidth]{figures/dynamic_curves.pdf}
    \vspace{5cm} % Placeholder space for the image
    \caption{SNDR and ENOB as a function of the input signal frequency and amplitude.}
    \label{fig:dynamic_curves}
\end{figure}

 

\section{Static Linearity Results}
\label{sec:static_linearity}

Static linearity metrics, specifically Differential Non-Linearity (DNL) and Integral Non-Linearity (INL), were evaluated to quantify the architectural and circuit-level deterministic mismatches. In an asynchronous Flash ADC, DNL and INL are heavily influenced by the random threshold variations of the comparator array and the parasitic resistance profile of the reference ladder. 

To counteract these non-idealities, a body bias trimming technique was proposed and implemented via a Verilog-A behavioral framework, as described in Chapter 5. To demonstrate the efficacy of this calibration mechanism, the static error profiles were extracted before and after the application of the optimal trimming vectors across multiple process corners. 

Figure~\ref{fig:dnl_inl_comparison} illustrates the untrimmed versus trimmed DNL and INL profiles across the converter's full-scale digital codes, highlighting the reduction in peak errors and the elimination of potential missing codes.

\begin{figure}[htbp]
    \centering
    % \includegraphics[width=0.9\textwidth]{figures/dnl_inl_comparison.pdf}
    \vspace{6cm} % Placeholder space for the image
    \caption{Measured DNL and INL error profiles before (top) and after (bottom) the application of the body bias trimming calibration.}
    \label{fig:dnl_inl_comparison}
\end{figure}

The summary of the peak and root-mean-square (RMS) static errors extracted from the simulation runs is presented in Table~\ref{tab:static_linearity_summary}.

\begin{table}[htbp]
    \centering
    \caption{Static linearity performance (DNL and INL) summary before and after calibration.}
    \label{tab:static_linearity_summary}
    \renewcommand{\arraystretch}{1.3}
    \begin{tabular}{lccccc}
        \toprule
        & \multicolumn{2}{c}{\textbf{Before Trimming}} & & \multicolumn{2}{c}{\textbf{After Trimming}} \\
        \cmidrule{2-3} \cmidrule{5-6}
        \textbf{Metric (LSB)} & \textbf{Peak} & \textbf{RMS} & & \textbf{Peak} & \textbf{RMS} \\
        \midrule
        DNL                   &               &              & &               &              \\
        INL                   &               &              & &               &              \\
        \bottomrule
    \end{tabular}
\end{table}

 

\section{Power Dissipation and Energy Efficiency}
\label{sec:power_efficiency}

A defining feature of the asynchronous Flash ADC is its signal-dependent power profile. In direct contrast to traditional synchronous architectures that exhibit a static power baseline due to continuous clock toggling, the dynamic power of the proposed converter is governed strictly by the frequency of threshold-crossing events. 

To empirically validate this adaptive scaling behavior, the total power consumption was monitored across an input signal frequency sweep spanning from DC ($f_{in} = 0$) up to the maximum tracking bandwidth limit. Figure~\ref{fig:power_scaling} explicitly demonstrates this linear dependency, proving that the system enters a near-zero dynamic consumption state when the input stimulus exhibits low activity or becomes idle.

\begin{figure}[htbp]
    \centering
    % \includegraphics[width=0.75\textwidth]{figures/power_scaling.pdf}
    \vspace{5cm} % Placeholder space for the image
    \caption{Total and dynamic power consumption as a function of the input signal frequency ($f_{in}$).}
    \label{fig:power_scaling}
\end{figure}

To provide a mathematically rigorous assessment of energy efficiency in this non-uniform context, an adapted Walden Figure of Merit ($FoM_{W,adapted}$) was utilized. This metric incorporates the effective average sampling rate ($f_{s,avg}$) dictated by the crossing density, ensuring a fair benchmark against synchronous designs:
\begin{equation}
    \mathrm{FoM}_{W,adapted} = \frac{P_{total}}{2^{\mathrm{ENOB}} \cdot f_{s,avg}}
\end{equation}

The extracted power breakdown and the resulting energy efficiency figures under both ideal and transient noise conditions are tabulated in Table~\ref{tab:power_metrics_empty}.

\begin{table}[htbp]
    \centering
    \caption{Power dissipation breakdown and energy efficiency summary.}
    \label{tab:power_metrics_empty}
    \renewcommand{\arraystretch}{1.4}
    \begin{tabular}{lcc}
        \toprule
        \textbf{Performance Parameter}           & \textbf{Value} & \textbf{Unit}         \\
        \midrule
        Static Power Consumption ($P_{static}$)  &                & $\mu$W                \\
        Dynamic Power (w/o Noise)                &                & $\mu$W                \\
        Dynamic Power (with Transient Noise)     &                & $\mu$W                \\
        Total Power Dissipation                  &                & $\mu$W                \\
        \midrule
        $\mathrm{FoM}_{W}$ (w/o Noise)           &                & fJ/conv.-step         \\
        $\mathrm{FoM}_{W}$ (with Transient Noise)&                & fJ/conv.-step         \\
        \bottomrule
    \end{tabular}
\end{table}

 

\section{Comparative Analysis and Discussion}
\label{sec:comparative_analysis}

To contextualize the performance of the developed design, a comprehensive comparative analysis is carried out. Table~\ref{tab:state_of_the_art} contrasts the core performance metrics of the baseline Synchronous Reference Flash ADC detailed in Chapter 4, the proposed Asynchronous Flash ADC developed in this work, and recent state-of-the-art architectures compiled from the literature.

\begin{table}[htbp]
    \centering
    \caption{Performance comparison with the synchronous baseline and state-of-the-art literature.}
    \label{tab:state_of_the_art}
    \renewcommand{\arraystretch}{1.3}
    \small
    \begin{tabular}{lccccc}
        \toprule
        \textbf{Parameter} & \textbf{Synchronous Baseline} & \textbf{This Work} & \textbf{Ref. [X]} & \textbf{Ref. [Y]} & \textbf{Ref. [Z]} \\
                           & \textbf{(Chapter 4)}          & \textbf{(Asynchronous)} & & & \\
        \midrule
        Technology (nm)    & 16                            & 16                 &                   &                   &                   \\
        Supply Voltage (V) & 0.8                           & 0.8                &                   &                   &                   \\
        Resolution (Bits)  &                               &                    &                   &                   &                   \\
        $f_s$ or $f_{s,avg}$ (MS/s) &                      &                    &                   &                   &                   \\
        SNDR (dB)          &                               &                    &                   &                   &                   \\
        ENOB (Bits)        &                               &                    &                   &                   &                   \\
        Power ($\mu$W)     &                               &                    &                   &                   &                   \\
        $\mathrm{FoM}_W$ (fJ/c.-step) &                    &                    &                   &                   &                   \\
        Architecture       & Synchronous Flash             & Asynchronous Flash &                   &                   &                   \\
        \bottomrule
    \end{tabular}
\end{table}

The comparative data underscores the trade-offs inherent to the asynchronous paradigm. While the synchronous baseline maintains a lower circuit complexity in its control block, it suffers from a rigid power envelope that compromises efficiency during periods of low input activity. In contrast, the proposed asynchronous implementation achieves superior energy scaling and a highly competitive Walden Figure of Merit, particularly after the activation of the body bias trimming network which successfully recovers the design's effective resolution under process variations.

